\scanpage{524}
\mspage{565}

$\overline{\text{\uline{III}}}$ \deleted{Let us consider the}

Let us consider the extension $\mathrm{Gl}(2,\mathbb Z)^\wedge$ of $\mathcal G_{0,3}^{\deleted{\unsure{+}}\wedge}$ by $\{\pm1\}$:

$(1)$
\diagram{\begin{tikzcd}[ampersand replacement=\&,column sep=1.4em,row sep=1.6em]
\&\&1\arrow[d]\&1\arrow[d]\&\\
1\arrow[r]\&\{\pm1\}\arrow[r,no head]\arrow[d,equal]\&
\mathrm{Sl}(2,\widehat{\mathbb Z})\arrow[r]\arrow[d]\&
\mathcal G_{0,3}^{+\wedge}\arrow[r]\arrow[d]\&1\\
1\arrow[r,no head]\&\{\pm1\}\arrow[r,no head]\&
\mathrm{Gl}(2,\widehat{\mathbb Z})\arrow[r,no head]\arrow[d]\&
\mathcal G_{0,3}^{\wedge}\arrow[r,no head]\arrow[d]\&1\\
\&\&\{\pm1\}\arrow[r]\arrow[d]\&\{\pm1\}\arrow[d]\&\\
\&\&1\&1\&
\end{tikzcd}}

We are interested in the group $\deleted{\unsure{M}}N_{1,1}^{\sim}$, defined \inserted{extension of $\deleted{\mathcal G_{0,3}^{\wedge}}M_{0,3}^{\sim}$ by $\{\pm1\}$,} as the inverse image of the extension $\mathrm{Gl}(2,\widehat{\mathbb Z})$ of $\mathrm{Gl}(2,\widehat{\mathbb Z})/\pm1$ by $\{\pm1\}$, via the canonical homomorphism
\[
(2)\qquad
M_{0,3}^{\sim}\deleted{\mathcal G_{0,3}^{\wedge}}
\longrightarrow\mathrm{Gl}(2,\widehat{\mathbb Z})'
=\mathrm{Gl}(2,\widehat{\mathbb Z})/\pm1
\]
(cf. \S45 \unsure{page}\ldots), so that we find a homomorphism of exact sequences,

\deleted{$\unsure{\tau_\infty(\epsilon_0)=\epsilon_0^{-1}}$}

$(3)$
\diagram{\begin{tikzcd}[ampersand replacement=\&,column sep=1.5em,row sep=2em]
1\arrow[r]\&\deleted{\unsure{M_{0,3}}}\{\pm1\}\arrow[r,no head]\arrow[d,equal]\&
\unsure{N_{1,1}^{\sim}|N_{0,3}^{\sim}}\arrow[r,no head]\arrow[d]\&M_{0,3}^{\sim}\arrow[r]\arrow[d]\&1\\
1\arrow[r]\&\{\pm1\}\arrow[r]\&\mathrm{Gl}(2,\widehat{\mathbb Z})\arrow[r]
\arrow[d,draw=none,"{\deleted{\unsure{\downarrow}}}"]\&
\mathrm{Gl}(2,\widehat{\mathbb Z})'\arrow[r,no head]
\arrow[d,draw=none,"{\deleted{\unsure{\downarrow}}}"]\&1\\
\&\&\deleted{\unsure{1}}\&\deleted{\unsure{1}}\&
\end{tikzcd}}

We find, \deleted{for $\unsure{\mathcal T_{1,1}^{\wedge}}$} \inserted{\unsure{a}} as the inverse image of $\deleted{\mathrm{Sl}(2,\mathbb Z)}\mathcal G_{0,3}^{+\wedge}\subset$ \deleted{in} $M_{0,3}^{\sim}$ in $N_{1,1}$, a $\mathrm{Gl}(2,\mathbb Z)^\wedge$, deduc-
\margin{\uline{N.B.} $\mathcal T_{1,1}\simeq\unsure{\mathrm{Gl}(2,\mathbb Z)_{\wedge}|\uline{\mathrm{Gl}(2,\mathbb Z)}}$.}
\mspage{523}


\scanpage{525}

ing a homomorphism
\[
(4)\qquad
\unsure{\mathrm{Sl}(2,\mathbb Z)^\wedge|\mathrm{Gl}(2,\mathbb Z)^\wedge}
\longrightarrow\deleted{\unsure{M}}N_{1,1}^{\sim}
\]
whose restriction to $\mathrm{Sl}(2,\mathbb Z)^\wedge$ fits \deleted{fits} into a canonical exact sequence

$(5)$
\diagram{\begin{tikzcd}[ampersand replacement=\&,column sep=1.8em]
1\arrow[r,no head]\&
\underbrace{\mathrm{Sl}(2,\mathbb Z)^{\unsure{\wedge}}}_{\mathcal T_{1,1}^{+\wedge}}\arrow[r]\&
N_{1,1}^{\sim}\arrow[r]\&
\underset{\substack{\Vert\\\Gamma_{0,3}^{\sim}}}{\Gamma_{\mathbb Q}^{\sim}}\arrow[r,no head]\&1
\end{tikzcd}}

which in turn fits into the diagram

$(6)$
\diagram{\begin{tikzcd}[ampersand replacement=\&,column sep=1.7em,row sep=2em]
1\arrow[r]\&\mathrm{Sl}(2,\mathbb Z)^\wedge\arrow[r]\arrow[d,"{\text{epimorphism}}"]\&
N_{1,1}^{\sim}\arrow[r,no head]\arrow[d,"{\text{epimorphism}}"]\&
\Gamma_{\mathbb Q}^{\sim}\arrow[r,no head]\arrow[d,"{\text{epimorphism}}","\chi"']\&1\\
1\arrow[r]\&\mathrm{Sl}(2,\widehat{\mathbb Z})\arrow[r,no head]\&
\mathrm{Gl}(2,\widehat{\mathbb Z})\arrow[r]\&\widehat{\mathbb Z}^{*}\arrow[r,no head]\&1
\end{tikzcd}}

We can write
\[
\begin{gathered}
(7)\qquad M_{0,3}^{\sim}\simeq
\overbrace{M_{0,3}^{\sim}(0)}^{\simeq\Gamma_{\mathbb Q}^{\sim}}
\cdot\overbrace{\mathcal G_{0,3}^{+\wedge}}^{\simeq\mathrm{Sl}(2,\mathbb Z)^{\wedge\prime}\simeq\mathrm{Sl}(2,\mathbb Z)^\wedge/\pm1}\\
\text{(semidirect)}
\end{gathered}
\]
in terms of \unsure{this} decomposition, and taking into account that we have a canonical homomorphism
\[
(8)\qquad M_{0,3}^{\sim}(0)\longrightarrow\mathrm{Gl}(2,\widehat{\mathbb Z})^{\unsure{\sharp|\prime}}
\]
\unsure{that|which} taking \unsure{the|there} \deleted{\unsure{homomorphism composed of}} \inserted{restriction of} (2) to $M_{0,3}^{\sim}(0)$, we find a canonical isomorphism
\mspage{524}

\scanpage{526}
\mspage{566}

\[
 (9)\qquad
 N_{1,1}^{\sim}\simeq
 \underbrace{N_{1,1}^{\sim}(0)}_{M_{0,3}^{\sim}(0)\simeq\Gamma_{\mathbb Q}^{\sim}}
 .\,\mathrm{Sl}(2,\mathbb Z)^\wedge .
\]

Thus we have an operation of
\inserted{$N_{1,1}^{\sim}(0)\simeq$}$\Gamma_{\mathbb Q}^{\sim}$
on $\mathrm{Sl}(2,\mathbb Z)^\wedge$
which I would like to make explicit, \deleted{We are going to}
\unsure{let us use|using} the \deleted{\unsure{isomorphism}} cartesian diagram

\[
(10)\qquad
\diagram{\begin{tikzcd}[ampersand replacement=\&,column sep=large,row sep=large]
\mathrm{Sl}(2,\unsure{\mathbb Z|\widehat{\mathbb Z}})^\wedge
 \arrow[r] \arrow[d]
 \arrow[dr,phantom,"\text{cartesian}"]
 \& \mathcal{G}_{0,3}^{+\wedge}
 \arrow[r,phantom,"\simeq"]
 \arrow[dr,"{\text{induced by (2)}}"']
 \& \mathrm{Sl}(2,\mathbb Z)^\wedge/\pm1
 \arrow[d,"{\unsure{q|q'}\ \text{canonical}}"]
 \\
\mathrm{Sl}(2,\widehat{\mathbb Z}) \arrow[rr]
 \& {} \& \mathrm{Sl}(2,\widehat{\mathbb Z})/\pm1
\end{tikzcd}}
\]

An element $\unsure{\widetilde{\nu}|\widetilde{\gamma}}$ of
$\Gamma_{\mathbb Q}^{\sim}$ corresponds to a pair
\deleted{$(\alpha,\beta)$ \unsure{of elements}}

\[
(11)\qquad
(\alpha,\beta)\in
\deleted{\unsure{\mathcal{G}_{0,3}}}
\underbrace{\widehat{\Pi}_{0,3}^{\deleted{\unsure{\wedge|\prime}}\prime}}_{
 \unsure{\{1,\tau_\infty\}}.\widehat{\Pi}_{0,3}}
\times\widehat{\Pi}_{0,3}
\]

satisfying \deleted{\unsure{of|to}} the loop equation

\[
(12)\qquad
\alpha\rho(\alpha)^{-1}=\beta\sigma(\beta)^{-1}\ell_1^\nu
\qquad (\nu\in\widehat{\mathbb Z})
\]

which operates on $\mathcal{G}_{0,3}^{+\wedge}\simeq
\mathrm{Sl}(2,\mathbb Z)^{\wedge\unsure{\prime}}$ by

\[
(13)\qquad
\left\{\begin{aligned}
u(\rho)&=\operatorname{int}(\alpha)\rho=\zeta\rho,
&\zeta&=\alpha\rho(\alpha)^{-1},\\
u(\sigma)&=\operatorname{int}(\beta)\sigma=\eta\sigma,
&\eta&=\beta\sigma(\beta)^{-1}
\end{aligned}\right.
\]

which is compatible with \deleted{an}\inserted{the} operation
\deleted{(\unsure{side} \illegible{2})}
\inserted{denoted $\unsure{\underline{u}|\underline{\alpha}}$}
on $\mathrm{Sl}(2,\widehat{\mathbb Z})$ induced
by $\operatorname{int}(\underline{u})$, where
$\underline{u}\in\mathrm{Gl}(2,\widehat{\mathbb Z})$ is the element

\mspage{525}

\scanpage{527}

defined \deleted{by} as the image of $u$ under (8), more
precisely

\[
(14)\qquad
\unsure{\widetilde{u}|\underline{u}}=
\unsure{i_{T_0}|i_{\Gamma_0}}(\mu),\qquad
\text{where }\mu=\chi(u)\text{ is the multiplier, and}
\]

\[
(15)\qquad
\left\{\begin{aligned}
\unsure{i_{T_0}|i_{\Gamma_0}} &:G_m\longrightarrow
\mathrm{Gl}(2,\widehat{\mathbb Z})\\
\mu&\longmapsto\begin{pmatrix}\mu&0\\0&1\end{pmatrix}.
\end{aligned}\right.
\]

We will have, more precisely

\[
(16)\qquad
\left\{\begin{aligned}
\underline{u}(\underline{\rho})&=\underline{\zeta}(\rho)
&\text{where }\underline{\zeta}&=\underline{\alpha}\rho(\underline{\alpha})^{-1}\\
\underline{u}(\underline{\sigma})&=\unsure{\underline{\eta}}(\sigma)
&\text{where }\underline{\eta}&=\varepsilon_{\unsure{\varphi|\rho}}
\underline{\beta}\sigma(\underline{\beta})^{-1}
\end{aligned}\right.
\]

\margin{the underlined ones denote the elements of
$\mathrm{Gl}(2,\widehat{\mathbb Z})$ ---
$\underline{\alpha}$ and $\underline{\beta}$
\unsure{are not determined} \unsure{except up to|except up to a} sign,
but $\underline{\zeta},\underline{\eta}$ are
\unsure{unambiguous|uniquely determined}.}

\deleted{$\varepsilon=\varepsilon_2(\mu)$}

with, as a reminder

\[
(17)\qquad
\varepsilon=\varepsilon_2(\mu)\in\pm1
\qquad
\left\{\begin{aligned}
\varepsilon&=+1&&\text{iff }\mu\equiv1\ (4)\\
\varepsilon&=-1&&\text{iff }\mu\equiv-1\ (4)
\end{aligned}\right.
\]

Of course, $\mathrm{Sl}(2,\mathbb Z)^\wedge$ is
\uline{generated} \inserted{\unsure{topologically}} by $\rho$ and $\sigma$,
and we identify $\pm1$ with the \unsure{central|center} subgroup
of $\mathrm{Sl}(2,\mathbb Z)$ hence of
$\mathrm{Sl}(2,\mathbb Z)^\wedge$, we will also write
$-g$ for the product of $g\in\mathrm{Sl}(2,\mathbb Z)^\wedge$
by $-1$. We will therefore have in $\mathrm{Sl}(2,\mathbb Z)^\wedge$

\[
(18)\qquad
\left\{\begin{aligned}
u(\rho)&=\deleted{\operatorname{int}\unsure{\zeta}(\rho)=}
\operatorname{int}(\widetilde{\alpha})(\rho)=\zeta\rho\\
u(\unsure{\sigma|\rho})&=
\varepsilon\operatorname{int}(\widetilde{\beta})(\sigma)=\eta\sigma
\end{aligned}\right.
\]
\unsure{and}

with
\[
(19)\qquad
\left\{\begin{aligned}
\zeta&=\widetilde{\alpha}\rho(\widetilde{\alpha})^{-1},\qquad
\eta=\varepsilon\widetilde{\beta}\sigma(\widetilde{\beta})^{-1}\\
&\varepsilon=\varepsilon_2(\mu)\in\{\pm1\}\subset
\mathrm{Sl}(2,\mathbb Z)^\wedge,
\end{aligned}\right.
\]

\mspage{526}

\scanpage{528}
\mspage{567}

where now $\widetilde{\alpha}$ and $\widetilde{\beta}$
denote arbitrary liftings of $\alpha$ and $\beta$ to
$\mathrm{Sl}(2,\mathbb Z)^\wedge$ --- they are determined
only up to sign, but
\inserted{$\operatorname{int}(\widetilde{\alpha}),
\operatorname{int}(\widetilde{\beta}),$}
$\zeta$ and $\eta$ are unambiguously determined by $u$.
\deleted{If \unsure{now}}
\inserted{Let us also note:}
\[
(20)\qquad u(\varepsilon_0)=\varepsilon_0^\mu
\quad\deleted{\text{\unsure{for}}}
\]

If now we have a \inserted{\unsure{surjective}} homomorphism
\[
(20)\qquad N_{1,1}^{\sim}\xrightarrow{\Psi}G
\]
inducing a homomorphism
\[
(21)\qquad
\mathrm{Sl}(2,\mathbb Z)^\wedge
\xrightarrow{\text{epimorphism}}
\underset{\substack{\Vert\\\operatorname{Im}\Psi}}{\mathcal{G}}
\subset G,
\]

the images of $\rho,\sigma$ in $\mathcal{G}$
(denoted by the same letters)
\unsure{satisfying|satisfy}
\[
(22)\qquad
\rho^6=\sigma^4=1,\qquad
\rho^3=\sigma^2\quad\text{and hence}
\]
\deleted{$[\rho^3,\sigma^2]=1$ \unsure{which suffices}}
\[
[\rho^3,\sigma]=[\sigma^2,\rho]=1
\quad\text{i.e. }\rho^3,\sigma^2
\]
\deleted{to \unsure{imply}}
are \uline{central}

\margin{thus in $\mathcal{G}$ there is a canonical
involutive central element, denoted $-1$ \unsure{. .}}

and we can, for $u\in N_{1,1}^{\sim}(0)$, corresponding to
\unsure{elements} $\widetilde{\alpha},\widetilde{\beta}$ of
$\mathrm{Sl}(2,\mathbb Z)^\wedge$
(determined modulo \uline{sign}), hence $\zeta,\eta$,
try to find an \inserted{\unsure{endomorphism|automorphism}}
$u_{\mathcal{G}}$ of $\mathcal{G}$,
satisfying the appropriate relations
\[
(23)\qquad
u_{\mathcal{G}}(\rho_{\mathcal{G}})
=\zeta_{\mathcal{G}}\rho_{\mathcal{G}},\qquad
u_{\mathcal{G}}(\sigma_{\mathcal{G}})
=\eta_{\mathcal{G}}\sigma_{\mathcal{G}},
\]
and bring into play \unsure{our|to the} $\widetilde{\alpha}$ and
$\widetilde{\beta}$. But we are

\mspage{527}

\scanpage{529}

obliged, since \deleted{\unsure{in general}} we do not
\inserted{\unsure{necessaril\textsuperscript{y}}}
$\widetilde{\alpha}\in\mathrm{Sl}(2,\mathbb Z)^\wedge$,
only $\widetilde{\alpha}\in\mathrm{Gl}(2,\mathbb Z)^\wedge$,
to bring in the element $\tau_\infty$ of
$\mathrm{Gl}(2,\mathbb Z)^\wedge$,
and \inserted{the action of} its image $\tau$ in $G$ on
$\mathcal{G}$. We will have, now taking

\[
(24)\qquad
\begin{aligned}
\overset{\deleted{\unsure{\mathcal{G}|\wedge}}}{\varepsilon_{0\mathcal{G}}}
&=\overset{\deleted{\unsure{\mathcal{G}|\wedge}}}{\sigma_{\mathcal{G}}}
\overset{\deleted{\unsure{\mathcal{G}|\wedge}}}{\rho_{\mathcal{G}}},
&\varepsilon_{1\mathcal{G}}&=\rho_{\mathcal{G}}\sigma_{\mathcal{G}}
\deleted{\unsure{=}}\,=\sigma_{\mathcal{G}}(\varepsilon_{0\mathcal{G}})
=\rho_{\mathcal{G}}(\varepsilon_{0\mathcal{G}})\\
&=-\sigma_{\mathcal{G}}^{-1}\rho_{\mathcal{G}}
&&=-\rho_{\mathcal{G}}\sigma_{\mathcal{G}}^{-1}
\end{aligned}
\]

\margin{which is not the usual convention
\unsure{up to} now\ldots}

the relations
\[
(25)\qquad
\tau(\sigma)=\sigma^{-1},\qquad
\tau(\rho)=\sigma(\rho^{-1})
\]

\margin{I drop the subscripts $\mathcal{G}$}

which imply, via (24) and (22)
\[
(26)\qquad
\tau(\varepsilon_0)=\varepsilon_0^{-1},\qquad
\tau(\varepsilon_1)=\varepsilon_1^{-1},
\]
and we introduce the group
\[
(27)\qquad
\deleted{\unsure{\{1,\tau\}}}
\mathcal{G}^{\unsure{\bullet|\prime}}
\overset{\mathrm{dfn}}{=}
\underset{\frac12\,\text{direct}}{\{1,\tau\}.\mathcal{G}},
\]
where $\{1,\tau\}$ operates on $\mathcal{G}$ by the formulas (25).

This said, we are in a position to paraphrase constructions
already made elsewhere.

\mspage{528}